\documentclass[12pt,openany]{book} % Default font size and left-justified equations

\usepackage[top=3cm,bottom=3cm,left=3.2cm,right=3.2cm,headsep=10pt,letterpaper]{geometry} % Page margins

\usepackage{xcolor} % Required for specifying colors by name
\definecolor{ocre}{RGB}{52,177,201} % Define the orange color used for highlighting throughout the book

% Font Settings
\usepackage{avant} % Use the Avantgarde font for headings
%\usepackage{times} % Use the Times font for headings
\usepackage{mathptmx} % Use the Adobe Times Roman as the default text font together with math symbols from the Sym­bol, Chancery and Com­puter Modern fonts

\usepackage{microtype} % Slightly tweak font spacing for aesthetics
\usepackage[utf8]{inputenc} % Required for including letters with accents
\usepackage[T1]{fontenc} % Use 8-bit encoding that has 256 glyphs

\usepackage[scaled]{helvet}
\renewcommand\familydefault{\sfdefault} 
\usepackage[T1]{fontenc}

% Bibliography
\usepackage[style=alphabetic,sorting=nyt,sortcites=true,autopunct=true,babel=hyphen,hyperref=true,abbreviate=false,backref=true,backend=biber]{biblatex}
\addbibresource{bibliography.bib} % BibTeX bibliography file
\defbibheading{bibempty}{}

%----------------------------------------------------------------------------------------
%	VARIOUS REQUIRED PACKAGES
%----------------------------------------------------------------------------------------

\usepackage{titlesec} % Allows customization of titles

\usepackage{graphicx} % Required for including pictures
\graphicspath{{Pictures/}} % Specifies the directory where pictures are stored

\usepackage{lipsum} % Inserts dummy text

\usepackage{tikz} % Required for drawing custom shapes

\usepackage[spanish]{babel}% language/hyphenation

\usepackage{enumitem} % Customize lists
\setlist{nolistsep} % Reduce spacing between bullet points and numbered lists

\usepackage{booktabs} % Required for nicer horizontal rules in tables

\usepackage{eso-pic} % Required for specifying an image background in the title page

%----------------------------------------------------------------------------------------
%	MAIN TABLE OF CONTENTS
%----------------------------------------------------------------------------------------

\usepackage{titletoc} % Required for manipulating the table of contents

\contentsmargin{0cm} % Removes the default margin
% Chapter text styling
\titlecontents{chapter}[1.25cm] % Indentation
{\addvspace{15pt}\large\sffamily\bfseries} % Spacing and font options for chapters
{\color{ocre!60}\contentslabel[\Large\thecontentslabel]{1.25cm}\color{ocre}} % Chapter number
{}  
{\color{ocre!60}\normalsize\sffamily\bfseries\;\titlerule*[.5pc]{.}\;\thecontentspage} % Page number
% Section text styling
\titlecontents{section}[1.25cm] % Indentation
{\addvspace{5pt}\sffamily\bfseries} % Spacing and font options for sections
{\contentslabel[\thecontentslabel]{1.25cm}} % Section number
{}
{\sffamily\hfill\color{black}\thecontentspage} % Page number
[]
% Subsection text styling
\titlecontents{subsection}[1.25cm] % Indentation
{\addvspace{1pt}\sffamily\small} % Spacing and font options for subsections
{\contentslabel[\thecontentslabel]{1.25cm}} % Subsection number
{}
{\sffamily\;\titlerule*[.5pc]{.}\;\thecontentspage} % Page number
[] 

%----------------------------------------------------------------------------------------
%	MINI TABLE OF CONTENTS IN CHAPTER HEADS
%----------------------------------------------------------------------------------------

% Section text styling
\titlecontents{lsection}[0em] % Indendating
{\footnotesize\sffamily} % Font settings
{}
{}
{}

% Subsection text styling
\titlecontents{lsubsection}[.5em] % Indentation
{\normalfont\footnotesize\sffamily} % Font settings
{}
{}
{}
 
%----------------------------------------------------------------------------------------
%	PAGE HEADERS
%----------------------------------------------------------------------------------------

\usepackage{fancyhdr} % Required for header and footer configuration

\pagestyle{fancy}
\renewcommand{\chaptermark}[1]{\markboth{\sffamily\normalsize\bfseries\chaptername\ \thechapter.\ #1}{}} % Chapter text font settings
\renewcommand{\sectionmark}[1]{\markright{\sffamily\normalsize\thesection\hspace{5pt}#1}{}} % Section text font settings
\fancyhf{} \fancyhead[LE,RO]{\sffamily\normalsize\thepage} % Font setting for the page number in the header
\fancyhead[LO]{\rightmark} % Print the nearest section name on the left side of odd pages
\fancyhead[RE]{\leftmark} % Print the current chapter name on the right side of even pages
\renewcommand{\headrulewidth}{0.5pt} % Width of the rule under the header
\addtolength{\headheight}{2.5pt} % Increase the spacing around the header slightly
\renewcommand{\footrulewidth}{0pt} % Removes the rule in the footer
\fancypagestyle{plain}{\fancyhead{}\renewcommand{\headrulewidth}{0pt}} % Style for when a plain pagestyle is specified

% Removes the header from odd empty pages at the end of chapters
\makeatletter
\renewcommand{\cleardoublepage}{
\clearpage\ifodd\c@page\else
\hbox{}
\vspace*{\fill}
\thispagestyle{empty}
\newpage
\fi}

%----------------------------------------------------------------------------------------
%	THEOREM STYLES
%----------------------------------------------------------------------------------------

\usepackage{amsmath,amsfonts,amssymb,amsthm} % For math equations, theorems, symbols, etc

\newcommand{\intoo}[2]{\mathopen{]}#1\,;#2\mathclose{[}}
\newcommand{\ud}{\mathop{\mathrm{{}d}}\mathopen{}}
\newcommand{\intff}[2]{\mathopen{[}#1\,;#2\mathclose{]}}
\newtheorem{notation}{Notation}[chapter]

%%%%%%%%%%%%%%%%%%%%%%%%%%%%%%%%%%%%%%%%%%%%%%%%%%%%%%%%%%%%%%%%%%%%%%%%%%%
%%%%%%%%%%%%%%%%%%%% dedicated to boxed/framed environements %%%%%%%%%%%%%%
%%%%%%%%%%%%%%%%%%%%%%%%%%%%%%%%%%%%%%%%%%%%%%%%%%%%%%%%%%%%%%%%%%%%%%%%%%%
\newtheoremstyle{ocrenumbox}% % Theorem style name
{0pt}% Space above
{0pt}% Space below
{\normalfont}% % Body font
{}% Indent amount
{\small\bf\sffamily\color{ocre}}% % Theorem head font
{\;}% Punctuation after theorem head
{0.25em}% Space after theorem head
{\small\sffamily\color{ocre}\thmname{#1}\nobreakspace\thmnumber{\@ifnotempty{#1}{}\@upn{#2}}% Theorem text (e.g. Theorem 2.1)
\thmnote{\nobreakspace\the\thm@notefont\sffamily\bfseries\color{black}---\nobreakspace#3.}} % Optional theorem note
\renewcommand{\qedsymbol}{$\blacksquare$}% Optional qed square

\newtheoremstyle{blacknumex}% Theorem style name
{5pt}% Space above
{5pt}% Space below
{\normalfont}% Body font
{} % Indent amount
{\small\bf\sffamily}% Theorem head font
{\;}% Punctuation after theorem head
{0.25em}% Space after theorem head
{\small\sffamily{\tiny\ensuremath{\blacksquare}}\nobreakspace\thmname{#1}\nobreakspace\thmnumber{\@ifnotempty{#1}{}\@upn{#2}}% Theorem text (e.g. Theorem 2.1)
\thmnote{\nobreakspace\the\thm@notefont\sffamily\bfseries---\nobreakspace#3.}}% Optional theorem note

\newtheoremstyle{blacknumbox} % Theorem style name
{0pt}% Space above
{0pt}% Space below
{\normalfont}% Body font
{}% Indent amount
{\small\bf\sffamily}% Theorem head font
{\;}% Punctuation after theorem head
{0.25em}% Space after theorem head
{\small\sffamily\thmname{#1}\nobreakspace\thmnumber{\@ifnotempty{#1}{}\@upn{#2}}% Theorem text (e.g. Theorem 2.1)
\thmnote{\nobreakspace\the\thm@notefont\sffamily\bfseries---\nobreakspace#3.}}% Optional theorem note

%%%%%%%%%%%%%%%%%%%%%%%%%%%%%%%%%%%%%%%%%%%%%%%%%%%%%%%%%%%%%%%%%%%%%%%%%%%
%%%%%%%%%%%%% dedicated to non-boxed/non-framed environements %%%%%%%%%%%%%
%%%%%%%%%%%%%%%%%%%%%%%%%%%%%%%%%%%%%%%%%%%%%%%%%%%%%%%%%%%%%%%%%%%%%%%%%%%
\newtheoremstyle{ocrenum}% % Theorem style name
{5pt}% Space above
{5pt}% Space below
{\normalfont}% % Body font
{}% Indent amount
{\small\bf\sffamily\color{ocre}}% % Theorem head font
{\;}% Punctuation after theorem head
{0.25em}% Space after theorem head
{\small\sffamily\color{ocre}\thmname{#1}\nobreakspace\thmnumber{\@ifnotempty{#1}{}\@upn{#2}}% Theorem text (e.g. Theorem 2.1)
\thmnote{\nobreakspace\the\thm@notefont\sffamily\bfseries\color{black}---\nobreakspace#3.}} % Optional theorem note
\renewcommand{\qedsymbol}{$\blacksquare$}% Optional qed square
\makeatother

% Defines the theorem text style for each type of theorem to one of the three styles above
\newcounter{dummy} 
\numberwithin{dummy}{section}
\theoremstyle{ocrenumbox}
\newtheorem{theoremeT}[dummy]{Theorem}
\newtheorem{problem}{Problem}[chapter]
\newtheorem{exerciseT}{Exercise}[chapter]
\theoremstyle{blacknumex}
\newtheorem{exampleT}{Example}[chapter]
\theoremstyle{blacknumbox}
\newtheorem{vocabulary}{Vocabulary}[chapter]
\newtheorem{definitionT}{Definition}[section]
\newtheorem{corollaryT}[dummy]{Corollary}
\theoremstyle{ocrenum}
\newtheorem{proposition}[dummy]{Proposition}

%----------------------------------------------------------------------------------------
%	DEFINITION OF COLORED BOXES
%----------------------------------------------------------------------------------------

\RequirePackage[framemethod=default]{mdframed} % Required for creating the theorem, definition, exercise and corollary boxes

% Theorem box
\newmdenv[skipabove=7pt,
skipbelow=7pt,
backgroundcolor=black!5,
linecolor=ocre,
innerleftmargin=5pt,
innerrightmargin=5pt,
innertopmargin=5pt,
leftmargin=0cm,
rightmargin=0cm,
innerbottommargin=5pt]{tBox}

% Exercise box	  
\newmdenv[skipabove=7pt,
skipbelow=7pt,
rightline=false,
leftline=true,
topline=false,
bottomline=false,
backgroundcolor=ocre!10,
linecolor=ocre,
innerleftmargin=5pt,
innerrightmargin=5pt,
innertopmargin=5pt,
innerbottommargin=5pt,
leftmargin=0cm,
rightmargin=0cm,
linewidth=4pt]{eBox}	

% Definition box
\newmdenv[skipabove=7pt,
skipbelow=7pt,
rightline=false,
leftline=true,
topline=false,
bottomline=false,
linecolor=ocre,
innerleftmargin=5pt,
innerrightmargin=5pt,
innertopmargin=0pt,
leftmargin=0cm,
rightmargin=0cm,
linewidth=4pt,
innerbottommargin=0pt]{dBox}	

% Corollary box
\newmdenv[skipabove=7pt,
skipbelow=7pt,
rightline=false,
leftline=true,
topline=false,
bottomline=false,
linecolor=gray,
backgroundcolor=black!5,
innerleftmargin=5pt,
innerrightmargin=5pt,
innertopmargin=5pt,
leftmargin=0cm,
rightmargin=0cm,
linewidth=4pt,
innerbottommargin=5pt]{cBox}

% Creates an environment for each type of theorem and assigns it a theorem text style from the "Theorem Styles" section above and a colored box from above
\newenvironment{theorem}{\begin{tBox}\begin{theoremeT}}{\end{theoremeT}\end{tBox}}
\newenvironment{exercise}{\begin{eBox}\begin{exerciseT}}{\hfill{\color{ocre}\tiny\ensuremath{\blacksquare}}\end{exerciseT}\end{eBox}}				  
\newenvironment{definition}{\begin{dBox}\begin{definitionT}}{\end{definitionT}\end{dBox}}	
\newenvironment{example}{\begin{exampleT}}{\hfill{\tiny\ensuremath{\blacksquare}}\end{exampleT}}		
\newenvironment{corollary}{\begin{cBox}\begin{corollaryT}}{\end{corollaryT}\end{cBox}}	

%----------------------------------------------------------------------------------------
%	REMARK ENVIRONMENT
%----------------------------------------------------------------------------------------

\newenvironment{remark}{\par\vspace{10pt}\small % Vertical white space above the remark and smaller font size
\begin{list}{}{
\leftmargin=35pt % Indentation on the left
\rightmargin=25pt}\item\ignorespaces % Indentation on the right
\makebox[-2.5pt]{\begin{tikzpicture}[overlay]
\node[draw=ocre!60,line width=1pt,circle,fill=ocre!25,font=\sffamily\bfseries,inner sep=2pt,outer sep=0pt] at (-15pt,0pt){\textcolor{ocre}{R}};\end{tikzpicture}} % Orange R in a circle
\advance\baselineskip -1pt}{\end{list}\vskip5pt} % Tighter line spacing and white space after remark

%----------------------------------------------------------------------------------------
%	SECTION NUMBERING IN THE MARGIN
%----------------------------------------------------------------------------------------

\makeatletter
\renewcommand{\@seccntformat}[1]{\llap{\textcolor{ocre}{\csname the#1\endcsname}\hspace{1em}}}                    
\renewcommand{\section}{\@startsection{section}{1}{\z@}
{-4ex \@plus -1ex \@minus -.4ex}
{1ex \@plus.2ex }
{\normalfont\large\sffamily\bfseries}}
\renewcommand{\subsection}{\@startsection {subsection}{2}{\z@}
{-3ex \@plus -0.1ex \@minus -.4ex}
{0.5ex \@plus.2ex }
{\normalfont\sffamily\bfseries}}
\renewcommand{\subsubsection}{\@startsection {subsubsection}{3}{\z@}
{-2ex \@plus -0.1ex \@minus -.2ex}
{.2ex \@plus.2ex }
{\normalfont\small\sffamily\bfseries}}                        
\renewcommand\paragraph{\@startsection{paragraph}{4}{\z@}
{-2ex \@plus-.2ex \@minus .2ex}
{.1ex}
{\normalfont\small\sffamily\bfseries}}

%----------------------------------------------------------------------------------------
%	HYPERLINKS IN THE DOCUMENTS
%----------------------------------------------------------------------------------------

% For an unclear reason, the package should be loaded now and not later
\usepackage{hyperref}
\hypersetup{hidelinks,backref=true,pagebackref=true,hyperindex=true,colorlinks=false,breaklinks=true,urlcolor= ocre,bookmarks=true,bookmarksopen=false,pdftitle={Title},pdfauthor={Author}}

%----------------------------------------------------------------------------------------
%	CHAPTER HEADINGS
%----------------------------------------------------------------------------------------

% The set-up below should be (sadly) manually adapted to the overall margin page septup controlled by the geometry package loaded in the main.tex document. It is possible to implement below the dimensions used in the goemetry package (top,bottom,left,right)... TO BE DONE

\newcommand{\thechapterimage}{}
\newcommand{\chapterimage}[1]{\renewcommand{\thechapterimage}{#1}}

% Numbered chapters with mini tableofcontents
\def\thechapter{\arabic{chapter}}
\def\@makechapterhead#1{
\thispagestyle{empty}
{\centering \normalfont\sffamily
\ifnum \c@secnumdepth >\m@ne
\if@mainmatter
\startcontents
\begin{tikzpicture}[remember picture,overlay]
\node at (current page.north west)
{\begin{tikzpicture}[remember picture,overlay]
\node[anchor=north west,inner sep=0pt] at (0,0) {\includegraphics[width=\paperwidth]{\thechapterimage}};
%%%%%%%%%%%%%%%%%%%%%%%%%%%%%%%%%%%%%%%%%%%%%%%%%%%%%%%%%%%%%%%%%%%%%%%%%%%%%%%%%%%%%
% Commenting the 3 lines below removes the small contents box in the chapter heading
%\fill[color=ocre!10!white,opacity=.6] (1cm,0) rectangle (8cm,-7cm);
%\node[anchor=north west] at (1.1cm,.35cm) {\parbox[t][8cm][t]{6.5cm}{\huge\bfseries\flushleft \printcontents{l}{1}{\setcounter{tocdepth}{2}}}};
\draw[anchor=west] (5cm,-9cm) node [rounded corners=20pt,fill=ocre!10!white,text opacity=1,draw=ocre,draw opacity=1,line width=1.5pt,fill opacity=.6,inner sep=12pt]{\huge\sffamily\bfseries\textcolor{black}{\thechapter. #1\strut\makebox[22cm]{}}};
%%%%%%%%%%%%%%%%%%%%%%%%%%%%%%%%%%%%%%%%%%%%%%%%%%%%%%%%%%%%%%%%%%%%%%%%%%%%%%%%%%%%%
\end{tikzpicture}};
\end{tikzpicture}}
\par\vspace*{230\p@}
\fi
\fi}

% Unnumbered chapters without mini tableofcontents (could be added though) 
\def\@makeschapterhead#1{
\thispagestyle{empty}
{\centering \normalfont\sffamily
\ifnum \c@secnumdepth >\m@ne
\if@mainmatter
\begin{tikzpicture}[remember picture,overlay]
\node at (current page.north west)
{\begin{tikzpicture}[remember picture,overlay]
\node[anchor=north west,inner sep=0pt] at (0,0) {\includegraphics[width=\paperwidth]{\thechapterimage}};
\draw[anchor=west] (5cm,-9cm) node [rounded corners=20pt,fill=ocre!10!white,fill opacity=.6,inner sep=12pt,text opacity=1,draw=ocre,draw opacity=1,line width=1.5pt]{\huge\sffamily\bfseries\textcolor{black}{#1\strut\makebox[22cm]{}}};
\end{tikzpicture}};
\end{tikzpicture}}
\par\vspace*{230\p@}
\fi
\fi
}
\makeatother
 % Insert the commands.tex file which contains the majority of the structure behind the template
\newcommand{\Cliente}{ENAEX S.A.}
\newcommand{\Producto}{EXPLOSIVO PX}
\newcommand{\Keywords}{['HOLLOW', 'EXPLOSIVE', 'CONIC', 'BOOSTER', 'SHAPED', 'CHARGE', 'HOLLOW', 'CHARGE']}
\newcommand{\Periodo}{1 DE JUNIO AL 30 DE JUNIO}

\begin{document}

%----------------------------------------------------------------------------------------
%	TITLE PAGE
%----------------------------------------------------------------------------------------

\begingroup
\thispagestyle{empty}
\AddToShipoutPicture*{\put(0,0){\includegraphics[scale=1.2]{cover1}}} % Image background
\vspace*{5cm}
\begin{flushleft}
\par\normalfont\fontsize{35}{35}\sffamily\selectfont
\textbf{Vigilancia}\\
\textbf{Tecnológica}\\
{\LARGE Reporte, Julio 2017}\par % Book title
\vspace*{1cm}
{\Huge Cliente: \Cliente}\par % Author name
\end{flushleft}
\endgroup

%----------------------------------------------------------------------------------------
%	COPYRIGHT PAGE
%----------------------------------------------------------------------------------------

\newpage
~\vfill
\thispagestyle{empty}

\noindent \textsc{Technology forecasting report}\\

\noindent \textsc{www.easypatents.cl}\\ % URL

\noindent Copyright \copyright 2017 , EasyPatents Technologies all rights reserved \\

\noindent \textit{\today} % Fecha en que se envia

%----------------------------------------------------------------------------------------
%	TABLE OF CONTENTS
%----------------------------------------------------------------------------------------

\chapterimage{head1.png} % Table of contents heading image

\pagestyle{empty} % No headers

\tableofcontents % Print the table of contents itself

%\cleardoublepage % Forces the first chapter to start on an odd page so it's on the right

\pagestyle{fancy} % Print headers again

%----------------------------------------------------------------------------------------
%	CHAPTER 1
%----------------------------------------------------------------------------------------
\chapterimage{head1.png} % Chapter heading image

\chapter{Introducción}
El presente Informe es una vigilancia de las solicitudes de patente y modelo de utilidad que se han presentado ante la WORLD INTELLECTUAL PROPERTY ORGANIZATION (WIPO) utilizando el TRATADO DE COOPERACIÓN EN MATERIA DE PATENTES (PCT) y con prioridad vigente para entrar a los siguientes países resaltados en azul. Lista completa al final del Informe con fecha de entrada en vigencia del Tratado.


%\chapterimage{head1.png} % Chapter heading image

\chapter{Metodología}
Una vez realizado el procedimiento PCT ante la WIPO, el solicitante tiene el derecho a la solicitud de exclusividad de comercialización, distribución, importación, fabricación y uso del producto o proceso en cada uno de los países suscritos al acuerdo en un plazo de 30 meses desde la fecha de solicitud PCT o prioridad.\\

A continuación se informan las solicitudes de patentes de invención y modelos de utilidad través del título, resumen, fecha de publicación, fecha de solicitud, inventores y solicitantes obtenidas en relación al producto \textsc{\Producto} asociados a las palabras clave \textsc{\Keywords} de \textsc{\Cliente} y que han sido publicados en el período correspondiente al plazo máximo de prioridad \textsc{\Periodo}.

\chapter{Resultados}
 A continuación se mostrarán los 10 resultados más 
    relevantes de la búsqueda, en caso de existir, el resto de los resultados pueden ser vistos en anexos.
 
 \vspace{1cm}1- \textbf{Número de publicación:} \href{https://worldwide.espacenet.com/publicationDetails/biblio?DB=EPODOC&II=0&ND=3&adjacent=true&locale=en_EP&FT=D&date=20160331&CC=WO&NR=2010088209A8&KC=A8#}{\textcolor{blue}{WO2010088209A8}}\\ 
\textbf{Titulo:} BOTTLE WITH DIRECTED POUR SPOUT\\ 
 
\textbf{Fecha de publicación:} 05 de Jan de 2010\\ 
\textbf{Fecha para ingreso a fases nacionales:} 05 de Jul de 2012\\ 
\textbf{Abstract:} described is a \colorbox{yellow}{bottle} having \colorbox{yellow}{bottle} central part having a \colorbox{yellow}{bottle} base, a handle, and a \colorbox{yellow}{bottle} top with a \colorbox{yellow}{bottle} filling fitment. a snap-on spout has a nozzle fitment with a forward directed pouring orifice and attached child-resistant nozzle. the snap-on spout is attached to the \colorbox{yellow}{bottle} central part with a non-circular snap-fit attachment \colorbox{yellow}{surface}.\\ 
 

 \vspace{1cm}2- \textbf{Número de publicación:} \href{https://worldwide.espacenet.com/publicationDetails/biblio?DB=EPODOC&II=0&ND=3&adjacent=true&locale=en_EP&FT=D&date=20160331&CC=WO&NR=1999014132A1&KC=A1#}{\textcolor{blue}{WO1999014132A1}}\\ 
\textbf{Titulo:} FORMED AND DECORATED SEAL\\ 
 
\textbf{Fecha de publicación:} 25 de Jan de 1999\\ 
\textbf{Fecha para ingreso a fases nacionales:} 25 de Jul de 2001\\ 
\textbf{Abstract:} a closure assembly (12, 14) for wine \colorbox{yellow}{bottle}s includes a \colorbox{yellow}{bottle} (10) having a cylindrical opening (18) in the neck (16), a funnel shaped \colorbox{yellow}{surface} (22) connected to the cylindrical opening and a flat planar rim (24) on the outside of the \colorbox{yellow}{bottle}. a secondary closure (14) in the shape of an inverted hat fits over the rim and engages a glass by means of a thermo\colorbox{yellow}{plastic} adhesive (36).\\ 
 

 \vspace{1cm}3- \textbf{Número de publicación:} \href{https://worldwide.espacenet.com/publicationDetails/biblio?DB=EPODOC&II=0&ND=3&adjacent=true&locale=en_EP&FT=D&date=20160331&CC=WO&NR=2008027029A3&KC=A3#}{\textcolor{blue}{WO2008027029A3}}\\ 
\textbf{Titulo:} TWO-PIECE CONTAINER SEAL AND METHOD OF MANUFACTURE\\ 
 
\textbf{Fecha de publicación:} 06 de Jan de 2008\\ 
\textbf{Fecha para ingreso a fases nacionales:} 06 de Jul de 2010\\ 
\textbf{Abstract:} an improved two-piece container seal suitable for sealing a container having a removable closure and for lining the closure comprises a flexible sealant sheet and a flexible liner sheet bonded thereto. the sealant sheet and liner sheet are of the same size and shape. the sealant sheet has a sealing \colorbox{yellow}{surface} and a first thermo\colorbox{yellow}{plastic} \colorbox{yellow}{surface}, and the liner sheet has an outer \colorbox{yellow}{surface} and a second thermo\colorbox{yellow}{plastic} \colorbox{yellow}{surface}. the first and second thermo\colorbox{yellow}{plastic} \colorbox{yellow}{surface}s are in opposed, congruent contact with each other, and the first thermo\colorbox{yellow}{plastic} \colorbox{yellow}{surface} is ultrasonically welded to the second thermo\colorbox{yellow}{plastic} \colorbox{yellow}{surface} by an array of spaced, frangible ultrasonic welds, forming an adhesive-free interface therebetween. a method of manufacturing a two-piece container seal and a sealed container comprising a two-piece container seal are also described.\\ 
 

 \vspace{1cm}4- \textbf{Número de publicación:} \href{https://worldwide.espacenet.com/publicationDetails/biblio?DB=EPODOC&II=0&ND=3&adjacent=true&locale=en_EP&FT=D&date=20160331&CC=WO&NR=2009018221A1&KC=A1#}{\textcolor{blue}{WO2009018221A1}}\\ 
\textbf{Titulo:} METHOD TO IMPROVE ADHESION OF A FORMED GASKET TO PLASTIC CLOSURES\\ 
 
\textbf{Fecha de publicación:} 05 de Jan de 2009\\ 
\textbf{Fecha para ingreso a fases nacionales:} 06 de Jul de 2011\\ 
\textbf{Abstract:} a method of attaching a gasket to a \colorbox{yellow}{plastic} \colorbox{yellow}{surface} includes providing a \colorbox{yellow}{plastic} material having a \colorbox{yellow}{surface}, applying heat to the \colorbox{yellow}{surface}, providing a liquid gasket material on the heated \colorbox{yellow}{surface}, and allowing the liquid gasket material and the \colorbox{yellow}{surface} to cool.\\ 
 

 \vspace{1cm}5- \textbf{Número de publicación:} \href{https://worldwide.espacenet.com/publicationDetails/biblio?DB=EPODOC&II=0&ND=3&adjacent=true&locale=en_EP&FT=D&date=20160331&CC=WO&NR=1997026198A1&KC=A1#}{\textcolor{blue}{WO1997026198A1}}\\ 
\textbf{Titulo:} CRYOGENIC FREEZING VIAL\\ 
 
\textbf{Fecha de publicación:} 24 de Jan de 1997\\ 
\textbf{Fecha para ingreso a fases nacionales:} 25 de Jul de 1999\\ 
\textbf{Abstract:} a cryogenic freezing vial (10). the vial (10) includes a hollow sample container (12) having exterior threads (26) seated inside a hollow outer container (16) having interior threads (30). the two containers define a surrounding volume (22) therebetween which operates as a thermal barrier to resist temperature change within the sample container (12). a cap (20) having a tubular wall (34) extending outwardly from a continuous capping \colorbox{yellow}{surface} (36) includes threads (42, 46) formed on both of the opposing sides (40, 44) of the tubular wall (34). the inner threaded side (40) of the tubular wall (34) can be screwed into the hollow sample container (12), and then the outer threaded side (44) of the tubular wall (34) can be screwed into the hollow outer container (16). in such manner the cap (20) is disposed in simultaneous threaded engagement between the two containers to intercouple and seal the containers from their surrounding environment.\\ 
 

 \vspace{1cm}6- \textbf{Número de publicación:} \href{https://worldwide.espacenet.com/publicationDetails/biblio?DB=EPODOC&II=0&ND=3&adjacent=true&locale=en_EP&FT=D&date=20160331&CC=WO&NR=2008095368A1&KC=A1#}{\textcolor{blue}{WO2008095368A1}}\\ 
\textbf{Titulo:} A DISPOSABLE AND BURGLARY-RESISTING DRINKING WATER BOTTLE\\ 
 
\textbf{Fecha de publicación:} 14 de Jan de 2008\\ 
\textbf{Fecha para ingreso a fases nacionales:} 14 de Jul de 2010\\ 
\textbf{Abstract:} a disposable and burglary-resisting drinking water \colorbox{yellow}{bottle}, comprises a \colorbox{yellow}{bottle} body, a body neck, a aluminum foil and a \colorbox{yellow}{bottle} cap, the \colorbox{yellow}{bottle} body is interacted and integrated with the \colorbox{yellow}{bottle} cap, the \colorbox{yellow}{bottle} neck is a part of the \colorbox{yellow}{bottle} body, the \colorbox{yellow}{bottle} cap can be connected with the \colorbox{yellow}{bottle} neck by thread, a lug and a locking band are set on the \colorbox{yellow}{bottle} neck, a center burglary-resisting piece is set on the \colorbox{yellow}{bottle} cap, a burglary-resisting ring is set on the bottom of the \colorbox{yellow}{bottle} cap, the sealing aluminum foil is set on the opening of the \colorbox{yellow}{bottle} neck.\\ 
 

 \vspace{1cm}7- \textbf{Número de publicación:} \href{https://worldwide.espacenet.com/publicationDetails/biblio?DB=EPODOC&II=0&ND=3&adjacent=true&locale=en_EP&FT=D&date=20160331&CC=WO&NR=2010072039A1&KC=A1#}{\textcolor{blue}{WO2010072039A1}}\\ 
\textbf{Titulo:} A WINE BOTTLE WITH AN OPENER\\ 
 
\textbf{Fecha de publicación:} 01 de Jan de 2010\\ 
\textbf{Fecha para ingreso a fases nacionales:} 01 de Jul de 2012\\ 
\textbf{Abstract:} a wine \colorbox{yellow}{bottle} with an opener includes a \colorbox{yellow}{bottle} cap (1), a \colorbox{yellow}{bottle} opening (2) and a \colorbox{yellow}{bottle} body (4). said \colorbox{yellow}{bottle} cap (1) is sleeved on the \colorbox{yellow}{bottle} opening (2) of the wine \colorbox{yellow}{bottle}, several external threads (3) are evenly provided on the external \colorbox{yellow}{surface} of said \colorbox{yellow}{bottle} opening (2), and internal threads matching with the external threads are formed on the internal \colorbox{yellow}{surface} of said \colorbox{yellow}{bottle} cap (1) due to the pressure on said \colorbox{yellow}{bottle} cap (1). a petal-like concave opener with a figure similar to the \colorbox{yellow}{bottle} cap is provided on the bottom of the \colorbox{yellow}{bottle}.\\ 
 

 \vspace{1cm}8- \textbf{Número de publicación:} \href{https://worldwide.espacenet.com/publicationDetails/biblio?DB=EPODOC&II=0&ND=3&adjacent=true&locale=en_EP&FT=D&date=20160331&CC=WO&NR=2003059768A1&KC=A1#}{\textcolor{blue}{WO2003059768A1}}\\ 
\textbf{Titulo:} PLASTIC CLOSING CAP WITH SEALING RING\\ 
 
\textbf{Fecha de publicación:} 24 de Jan de 2003\\ 
\textbf{Fecha para ingreso a fases nacionales:} 24 de Jul de 2005\\ 
\textbf{Abstract:} a closing cap moulded in one piece from a \colorbox{yellow}{plastic}s material for sealing a container with a threaded neck has a top plate (3), an annular skirt (4) with an internal screw thread (5) adapted to cooperate with the threaded container neck (2), and an integrally moulded deformable annular sealing ring (6) extending downwardly from an inner \colorbox{yellow}{surface} of said top plate (3) and spaced radially inwardly of the annular skirt (4), wherein when the closing cap (1) is screwed onto the container neck (2), the sealing ring (6) is deformed outwardly and is fully conformed to the outer \colorbox{yellow}{surface} of the container neck (2) without being clamped between the container neck (2) and the annular skirt (4) and without being pressed against the annular skirt (4).\\ 
 

 \vspace{1cm}9- \textbf{Número de publicación:} \href{https://worldwide.espacenet.com/publicationDetails/biblio?DB=EPODOC&II=0&ND=3&adjacent=true&locale=en_EP&FT=D&date=20160331&CC=WO&NR=1988003115A1&KC=A1#}{\textcolor{blue}{WO1988003115A1}}\\ 
\textbf{Titulo:} PLASTIC CLOSURE WITH UNITARILY MOLDED, FOAMED SEALING LAYER\\ 
 
\textbf{Fecha de publicación:} 05 de Jan de 1988\\ 
\textbf{Fecha para ingreso a fases nacionales:} 05 de Jul de 1990\\ 
\textbf{Abstract:} a molded polymeric container closure (10) comprising at least one substantially unfoamed polymer (14a) layer and an integrally molded foamed layer (14b) of the same polymer.\\ 
 

 \vspace{1cm}10- \textbf{Número de publicación:} \href{https://worldwide.espacenet.com/publicationDetails/biblio?DB=EPODOC&II=0&ND=3&adjacent=true&locale=en_EP&FT=D&date=20160331&CC=WO&NR=1993003870A1&KC=A1#}{\textcolor{blue}{WO1993003870A1}}\\ 
\textbf{Titulo:} LEAD-FREE CAPSULE FOR WINE BOTTLES AND METHOD OF MAKING\\ 
 
\textbf{Fecha de publicación:} 04 de Jan de 1993\\ 
\textbf{Fecha para ingreso a fases nacionales:} 05 de Jul de 1995\\ 
\textbf{Abstract:} a lead-free capsule (26) for wine \colorbox{yellow}{bottle}s and the like and a method for producing the capsule to have similar characteristics as traditional lead and lead-tin capsules include the selection of a laminate (10) of metallic foil (16) and \colorbox{yellow}{plastic} sheet (14) having a soft metallic \colorbox{yellow}{surface} look and feel, the laminate having characteristics of flexibility, ductility and adhesion permitting the laminate to be drawn to form a continuous, seamless capsule (26) configuration while maintaining laminate integrity.\\ 
 

 \vspace{1cm}11- \textbf{Número de publicación:} \href{https://worldwide.espacenet.com/publicationDetails/biblio?DB=EPODOC&II=0&ND=3&adjacent=true&locale=en_EP&FT=D&date=20160331&CC=WO&NR=2014139352A1&KC=A1#}{\textcolor{blue}{WO2014139352A1}}\\ 
\textbf{Titulo:} SAFETY ANTI-COUNTERFEITING WINE BOTTLE AND STRUCTURE OF BOTTLE CAP\\ 
 
\textbf{Fecha de publicación:} 18 de Jan de 2014\\ 
\textbf{Fecha para ingreso a fases nacionales:} 18 de Jul de 2016\\ 
\textbf{Abstract:} the present invention relates to a safety anti-counterfeiting wine \colorbox{yellow}{bottle} and a structure of a \colorbox{yellow}{bottle} cap. the \colorbox{yellow}{bottle} comprises a \colorbox{yellow}{bottle} body, a \colorbox{yellow}{bottle} cap and an anti-counterfeiting ring, wherein the anti-counterfeiting ring is a metallic anti-counterfeiting ring which is stuck to the outer wall of a \colorbox{yellow}{bottle} mouth of the wine \colorbox{yellow}{bottle}, the \colorbox{yellow}{bottle} mouth of the wine \colorbox{yellow}{bottle} is covered by the \colorbox{yellow}{bottle} cap, the \colorbox{yellow}{bottle} cap is a metallic \colorbox{yellow}{bottle} cap, a limiting convex rib which extends inwards is arranged on the top of the metallic anti-counterfeiting ring, a skirt edge which extends outwards is arranged on the bottom of the metallic \colorbox{yellow}{bottle} cap, the outer diameter of the skirt edge is greater than the inner diameter of the limiting convex rib, the skirt edge is arranged in the metallic anti-counterfeiting ring, a sealing ring is also arranged on the outer wall of the \colorbox{yellow}{bottle} mouth of the \colorbox{yellow}{bottle} body, and is used for sealing the inner wall of the \colorbox{yellow}{bottle} cap and the outer wall of the \colorbox{yellow}{bottle} mouth, a circle of easily-broken grooves are arranged on the side wall of the \colorbox{yellow}{bottle} cap, and a rotary clamping position for rotating the \colorbox{yellow}{bottle} cap is also arranged on the top of the \colorbox{yellow}{bottle} cap. in the present invention, by adopting the metallic sealing structure of the \colorbox{yellow}{bottle} cap, not only the sealing performance of the wine \colorbox{yellow}{bottle} is ensured, but also the problem of \colorbox{yellow}{plastic}izer of the existing sealing structure of the wine \colorbox{yellow}{bottle} is solved.\\ 
 


\end{document}
